\subsection{Estructura de la investigación}

La organización de esta ICR retoma la propuesta de estructurar la investigación a partir de una relación explícita entre problemática, objeto de estudio, unidad de análisis, preguntas, objetivos, procedimiento, límites y producto \citep{guerrerobaca2025propuesta}. Esta relación permite mantener el trabajo acotado a la evidencia disponible y evita que el diagnóstico prometa alcances que no forman parte del corpus.

\begin{table}[ht]
  \centering
  \caption{Estructura operativa de la investigación}
  \vspace{0.25cm}
  \label{tab:estructura_investigacion}
  \begin{tabular}{@{}L{0.28\linewidth}L{0.62\linewidth}@{}}
    \toprule
    \textbf{Elemento} & \textbf{Definición en esta ICR} \\
    \midrule
    Problemática & Transformación, abandono y baja documentación de construcciones tradicionales hñähñu en dos comunidades rurales del Valle del Mezquital. \\
    Objeto de estudio & Permanencia y transformación de sistemas constructivos tradicionales vinculados con materiales locales, uso cotidiano y memoria oral. \\
    Unidad de análisis & Cuatro casos seleccionados de construcciones tradicionales, codificados como C-01, C-02, C-03 y C-04, dentro de un registro más amplio. \\
    Categorías de análisis & Sistema constructivo, materiales, uso, estado físico, cambios y permanencias, relato asociado, ubicación cartográfica y manejo de privacidad. \\
    Evidencia empírica & Fotografías, geolocalizaciones, recorridos, relatos hablados y notas de observación disponibles. \\
    Procedimiento & Ordenamiento del corpus, elaboración de fichas, análisis cartográfico, lectura de relatos y síntesis comparativa. \\
    Producto & Fichas descriptivas, mapas generalizados, cuadro comparativo y criterios básicos de documentación y conservación preventiva. \\
    Alcance & Diagnóstico descriptivo y cartográfico; no incluye levantamientos arquitectónicos medidos, encuestas socioeconómicas ni propuesta de restauración. \\
    \bottomrule
  \end{tabular}
\end{table}

Esta estructura también define la escala del documento: los resultados no pretenden representar la totalidad de la arquitectura vernácula hñähñu, sino organizar con rigor cuatro casos seleccionados por su valor morfológico y material para reconocer diversidad constructiva, transformaciones recientes y posibilidades de documentación posterior.
